
\section{Limitations}

Several limitations of the present work should be kept in mind when interpreting
the results of Chapter~\ref{ch:results} and Chapter~\ref{ch:results_1LL}.

\paragraph{Finite-size effects.} All results are obtained by exact diagonalization
on a single system size (a $3 \times 4$ torus for the phase diagrams, with a
separate tilted geometry for the momentum-occupation calculations of \S\ref{sec:third_third_intermediate_phases}
onward). As discussed in \S\ref{sec:recognition}, a rigorous identification of topological order
requires finite-size scaling --- confirming that the many-body gap remains open and
the ground-state degeneracy persists as the system size increases --- which was not
performed here due to computational cost. The phase boundaries reported in
Chapter~\ref{ch:results} and Chapter~\ref{ch:results_1LL} should therefore be understood as consistent with, rather than a
definitive proof of, the identified phases; in particular, the precise location of
phase boundaries (e.g.\ $\alpha_c$, $W_c$) is expected to shift with system size,
even where the qualitative phase structure is robust.

\paragraph{Momentum grid choice.} Related to the above, the untilted rectangular
momentum grid used by default (\S\ref{sec:tilted_bc}) is known to bias the competition between
topological and charge-density-wave-like states relative to more isotropic tilted
grids~\cite{lauchli2013hierarchy}. While this choice was made for simplicity and
consistency across the phase diagrams, it means the reported phase boundaries may
not be fully representative of the thermodynamic-limit behavior.

\paragraph{Treatment of the $q=0$ interaction component.} As noted in \S\ref{sec:interactions}, the
``background-subtracted'' variant of the interaction removes the full
$q \to 0$ component of $\tilde{U}_{\tau\tau'}(q)$ for all $\tau, \tau'$, rather than
only the part that couples to the total charge density. This is a simplification of
the standard neutralizing-background convention, and the $\alpha$-dependence of the
discarded component is in principle physical. The qualitative difference between the
two phase diagrams shown throughout Chapter~4 (with and without this subtraction)
suggests this choice is not negligible, and a more careful treatment is left for
future work.

\paragraph{Model scope.} The model of Chapter~2 captures the two-valley
Landau-level physics of Moir\'e TMDs at a single magic angle, projected onto a
single (Landau-level-like) band per valley, and neglects Landau-level mixing. Real Moir\'e TMD flat bands deviate from
ideal Landau levels in ways not captured by this construction (e.g.\ non-uniform
Berry curvature, deviations of the form factor from the ideal Gaussian form), and a
quantitative comparison to specific materials (e.g.\ MoTe$_2$ at
$\theta \approx 1.89^\circ$) would require using the actual continuum-model form
factors rather than the idealized Landau-level ones used here.

\paragraph{Unresolved intermediate phases.} Several regions of the phase diagrams
-- most notably the narrow 3-fold-degenerate region near the FTI--gapless boundary
at $\nu = 2/3$ (\S4.3.3), the candidate IVC regions inferred from flat momentum
occupation at $\nu = 2/3$ and $\nu = 1$ (\S4.3.3, \S4.4), and the
$|\mathrm{CFL}\rangle \otimes |\mathrm{CFL}\rangle$--$|\mathrm{FL}\rangle \otimes
|\mathrm{FL}\rangle$ transition at $\nu = 1$ (\S4.4) -- were identified only
indirectly, through the degeneracy/spectral-flow/occupation diagnostics of \S3.4,
rather than through direct comparison with model wavefunctions. A more conclusive
identification of these states would require computing their overlap with
candidate trial wavefunctions (e.g.\ the Halperin $(1,1,2)$ state), which was
beyond the scope of this thesis.


% ------------------------------------------------------------
% 4. CONCLUSIONS
% Final chapter (or final section of Ch. 5, as preferred).
% ------------------------------------------------------------

\section{Conclusions}

This thesis studied a two-valley Landau-level model designed to capture the
essential physics of correlated and topological phases in Moir\'e transition metal
dichalcogenides, using exact diagonalization on finite tori. At filling
$\nu = 2/3$, we found clear numerical evidence --- via ground-state degeneracy and
spectral flow --- for a valley-polarized fractional quantum Hall
state at small inter-valley coupling $\alpha$ and small bandwidth $W$, and for a
valley-unpolarized fractional topological insulator built from two time-reversed
Laughlin copies at large $\alpha$, consistent with theoretical expectations. Both
phases were shown to survive a finite bandwidth and finite screening length, with
stability boundaries that depend systematically on $d_{\mathrm{SC}}$ and
$l_{\mathrm{RK}}$. At $\nu = 1$ and $\nu = 1/2$, we recovered the expected
valley-polarized integer quantum Hall state and decoupled composite-Fermi-liquid
pairs, respectively, together with signatures of an intervening
inter-valley-coherent regime at large $|W|$.

An analytic Hartree--Fock treatment of the valley-depolarization transition,
developed for both the $\alpha$-driven and $W$-driven cases, reproduced the correct
qualitative trends of the numerical phase diagrams -- including the direction of
the phase boundaries with $d_{\mathrm{SC}}$ and $l_{\mathrm{RK}}$ -- while
systematically overestimating the critical $\alpha_c$ and $W_c$ by up to an order
of magnitude, a discrepancy attributed largely to the circular-Fermi-surface
approximation used in the mean-field derivation.

Taken together, these results support the picture of Moir\'e TMDs as a
magnetic-field-free platform for realizing both fractional Chern insulator and
fractional topological insulator physics, and provide a systematic account of how
this physics competes with bandwidth, inter-valley coupling, and dielectric
screening within a single, tunable model. The main open questions raised by this
work --- the precise nature of the intermediate phases identified only indirectly in
\S4.3--4.4, the role of finite-size effects, and the extension to non-Abelian
Moore--Read-like states in the first excited Landau level --- are left for future work.